\documentclass[Total.tex]{subfiles}

\begin{document}
	\chapter{Valuation Rings}
	\begin{Def}
		Valuation rings:
		\begin{itemize}
			\item Let $K$ be a field with $A,B\subseteq K$ local rings. We say that $B$ \textbf{dominates} $A$ if $A\subseteq B$ and $m_A=A\cap m_B$;
			\item Let $A$ be a ring. We say $A$ is a \textbf{valuation ring} if $A$ is a local domain and if $A$ is maximal for the relation of domination among local rings contained in the fraction field of $A$;
			\item Let $A$ be a valuation ring with fraction field $K$. If $R\subseteq K$ is a subring of $K$, then we say $A$ is \textbf{centered} on $R$ if $R\subseteq A$.
		\end{itemize}
	\end{Def}
	
	\begin{lemma}
		(Existence of Valuation Ring) Let $K$ be a field and $A\subseteq K$ be a local ring. Then there exists a valuation ring with faction field $K$ dominating $A$.
	\end{lemma}
	\begin{proof}
		It is suffice to show: For collection of totally ordered collection of local subrings 
	\end{proof}
	
	\begin{lemma}
		Let $A$ be a valuation ring, then $A$ is a normal domain.
	\end{lemma}
	\begin{proof}
		内容...
	\end{proof}
	
	\begin{lemma}
		Let $A$ be a valuation ring with maximal ideal $m$ and fraction field $K$, and $x\in K$. Then, either $x\in A$ or $x^{-1}\in A$ or both.
	\end{lemma}
	\begin{proof}
		内容...
	\end{proof}
	\begin{lemma}
		内容...
	\end{lemma}
	\begin{proof}
		内容...
	\end{proof}
		\section{Valuation}
			\begin{lemma}
				Let $A$ be a valuation ring with field of fraction $K$. The set
				\begin{gather*}
					内容...
				\end{gather*}
			\end{lemma}
			\begin{proof}
				内容...
			\end{proof}
			\begin{Def}
				Let $A$ be a valuation ring:
				\begin{itemize}
					\item The totally ordered abelian group $(\Gamma,\geq)$ above is called \textbf{the value group} of valuation ring $A$;
					\item Map $v:A-\{0\}\rightarrow \Gamma$ and $v:K^*\rightarrow \Gamma$ is called valuation associated to $A$;
					\item The valuation ring $A$ is called \textbf{discrete valuation ring} if $\Gamma\cong \ZZ$.
				\end{itemize}
			\end{Def}
			
			\begin{lemma}
				内容...
			\end{lemma}
			\begin{proof}
				内容...
			\end{proof}
	\chapter{Valuation Theory}
	\chapter{Ultrametric Space}

\end{document}